% Template for the validation run of SMOS L2 against ISMN in situ data.
% ISMN is dataset 0, SMOS L2 dataset 1 in the netcdf results.
\def\runidx{\detokenize{$<ContentVars['report_run_index']>$}}
\def\runpair{0-\detokenize{$<ContentVars['NetcdfMetaVars']['val_dc_dataset0']>$}_and_1-\detokenize{$<ContentVars['NetcdfMetaVars']['val_dc_dataset1']>$}}

\subsection{Validation run \runidx{} -- \texorpdfstring{\acs{ISMN}}{ISMN} (\detokenize{$<ContentVars['ConfigVars']['DS1']['version']>$})}

\subsubsection{Validation settings}

The following validation settings were used for the comparison of \ac{SMOS} \ac{L2} against \ac{ISMN} in situ \ac{FRMs}.
For a full list of the settings used, see the results page of this validation run, linked in Table \ref{tab:val_run_list}.

\begin{itemize}
    \item Spatial reference (spatial matching): \textbf{\detokenize{$<ContentVars['ConfigVars']['spatial_reference']>$}}
    \item Temporal reference (temporal matching): \textbf{\detokenize{$<ContentVars['ConfigVars']['temporal_reference']>$}}
    \item Temporal matching window: \textbf{\detokenize{$<ContentVars['ConfigVars']['temporal_matching']>$} hours}
    \item Scaling reference: \textbf{\detokenize{$<ContentVars['ConfigVars']['scaling_reference']>$}}
    \item Validation metrics \detokenize{$<'\\textbf{are}' if ContentVars['ConfigVars']['anomalies_method']!='none' else '\\textbf{are not}'>$} based on anomalies.
\end{itemize}

\subsubsection{Results}

\noindent \ac{SMOS} \ac{L2} was compared against \textbf{\detokenize{$<ContentVars['NetcdfVars']['n_gpis']>$}} in situ time series.
Note that some of these sensors can be collocated, as a single monitoring site usually operates multiple sensors.
The nearest \ac{SMOS} grid point is found for each \ac{ISMN} location (nearest-neighbour search with \detokenize{$<ContentVars['ConfigVars']['spatial_reference']>$} as the reference); all validation results are therefore computed at the \ac{ISMN} locations.
\\[1em]
\noindent Validation metrics could be computed for \textbf{\detokenize{$<ContentVars['NetcdfVars']['n_ok'][0]>$}} sensors (\detokenize{$<ContentVars['NetcdfVars']['percent_ok'][0]>$}\,\%).
For the remaining \detokenize{$<ContentVars['NetcdfVars']['n_failed'][0]>$} sensors, no metrics were computed (Fig. \ref{fig:status\runidx}), mostly because too few valid observation pairs remained after filtering and temporal matching.
This typically affects stations with frozen or snow-covered soils (which are masked in \ac{SMOS} \ac{L2}), and networks that did not provide recent data.
Table \ref{tab:networks} breaks the sample down by network.

\begin{table}[htb!]
    \centering
    \small
    \caption{\ac{ISMN} networks in the validation run: number of tested sensors, number of sensors with valid results, and median Pearson's R and \ac{ubRMSD} (in m\textsuperscript{3}m\textsuperscript{-3}) of the valid sensors. References for each network are listed in Table \ref{tab:ismn_refs}.}
    \label{tab:networks}
    \begin{tabular}{@{}lrrrr@{}}
        \toprule
        \textbf{Network} & \textbf{Sensors} & \textbf{Valid} & \textbf{Median R} & \textbf{Median ubRMSD} \\
        \midrule
        \detokenize{$<' \\\\ '.join([n['network'].replace('_', '\\_') + ' & ' + str(n['n_sensors']) + ' & ' + str(n.get('n_ok', '--')) + ' & ' + ('--' if n.get('median_R') is None else format(n['median_R'], '.2f')) + ' & ' + ('--' if n.get('median_urmsd') is None else format(n['median_urmsd'], '.3f')) for n in ContentVars['NetcdfVars']['networks']]) + ' \\\\'>$}
        \bottomrule
    \end{tabular}
\end{table}

\noindent Table \ref{tab:sumstats\runidx} lists summary statistics across all valid sensors: the number of temporally matched observations used for the metric computation ("\#Obs."), Pearson's R, the bias (difference in means) and the \ac{ubRMSD}.
Additional metrics are available in the validation results \href{\detokenize{$<ContentVars['NetcdfMetaVars']['url']>$}}{online}.

\begin{table}[htb!]
    \centering
    \caption{Summary statistics of \ac{SMOS} \ac{L2} against \ac{ISMN} across all valid sensors.}
    \label{tab:sumstats\runidx}
    \begin{tabular}{ l r r r r }
        \toprule
        \textbf{Metric} & \textbf{\#Obs.} & \textbf{Pearson's R} & \textbf{Bias}  & \textbf{ubRMSD}  \\
        \midrule
        \textbf{Mean}   & \detokenize{$<ContentVars['SummaryVars']['MEAN_N. OBS.']>$}   & \detokenize{$<ContentVars['SummaryVars']["MEAN_PEARSON'S R"]>$}     & \detokenize{$<ContentVars['SummaryVars']['MEAN_BIAS']>$}   & \detokenize{$<ContentVars['SummaryVars']['MEAN_UNBIASED ROOT-MEAN-SQUARE DEVIATION']>$}   \\
        \textbf{Median}   & \detokenize{$<ContentVars['SummaryVars']['MEDIAN_N. OBS.']>$}   & \detokenize{$<ContentVars['SummaryVars']["MEDIAN_PEARSON'S R"]>$}     & \detokenize{$<ContentVars['SummaryVars']['MEDIAN_BIAS']>$}   & \detokenize{$<ContentVars['SummaryVars']['MEDIAN_UNBIASED ROOT-MEAN-SQUARE DEVIATION']>$}   \\
        \textbf{Interquartile range}   & \detokenize{$<ContentVars['SummaryVars']['IQ RANGE_N. OBS.']>$}   & \detokenize{$<ContentVars['SummaryVars']["IQ RANGE_PEARSON'S R"]>$}     & \detokenize{$<ContentVars['SummaryVars']['IQ RANGE_BIAS']>$}   & \detokenize{$<ContentVars['SummaryVars']['IQ RANGE_UNBIASED ROOT-MEAN-SQUARE DEVIATION']>$}   \\
        \bottomrule
    \end{tabular}
\end{table}

\begin{figure}[ht!]
\centering
\begin{minipage}[c]{0.34\textwidth}
    \centering
    \plotfig{./qa4sm_graphics/bulk_barplot_status_\runpair.png}
    \subcaption{} % (a)
\end{minipage}%
\hfill
\begin{minipage}[c]{0.64\textwidth}
    \centering
    \mapfig{./qa4sm_graphics/bulk_overview_\runpair_status.png}
    \subcaption{} % (b)
\end{minipage}
\caption{Number of successfully validated sensors and errors (a), and their location (b) in the validation run against \ac{ISMN}.}
\label{fig:status\runidx}
\end{figure}

\noindent Validation metrics are computed from temporally matched observations of \ac{SMOS} \ac{L2} and \ac{ISMN}.
A \detokenize{$<ContentVars['ConfigVars']['temporal_matching']>$}-hour temporal matching window is applied to each \detokenize{$<ContentVars['ConfigVars']['temporal_reference']>$} observation to find the closest data point in the in situ time series.
This minimises the potential (temporal) lag between the time series and the impact of (precipitation) events that occurred shortly after a satellite overpass.
The resulting sample size, which applies to all computed metrics, is shown in Fig. \ref{fig:nobs\runidx}.

\begin{figure}[ht!]
\centering
\begin{minipage}[c]{0.34\textwidth}
    \centering
    \plotfig{./qa4sm_graphics/bulk_boxplot_n_obs.png}
    \subcaption{} % (a)
\end{minipage}%
\hfill
\begin{minipage}[c]{0.64\textwidth}
    \centering
    \mapfig{./qa4sm_graphics/bulk_overview_n_obs.png}
    \subcaption{} % (b)
\end{minipage}
\caption{Distribution of the validation sample size (a) and sample size per location (b) in the validation run against \ac{ISMN}.}
\label{fig:nobs\runidx}
\end{figure}

\noindent Fig. \ref{fig:Rbox\runidx} shows Pearson's R between \ac{SMOS} \ac{L2} and \ac{ISMN} at all valid sensors.
The \ac{ubRMSD} is shown in Fig. \ref{fig:urmsd\runidx}, and Fig. \ref{fig:Rmeta\runidx} splits both metrics by the land cover class of the in situ sites.

\begin{figure}[ht!]
\centering
\begin{minipage}[c]{0.34\textwidth}
    \centering
    \plotfig{./qa4sm_graphics/bulk_boxplot_R.png}
    \subcaption{} % (a)
\end{minipage}%
\hfill
\begin{minipage}[c]{0.64\textwidth}
    \centering
    \mapfig{./qa4sm_graphics/bulk_overview_\runpair_R.png}
    \subcaption{} % (b)
\end{minipage}
\caption{Distribution of Pearson's R across all valid sensors (a) and spatial overview of Pearson's R (b) between \ac{SMOS} \ac{L2} and \ac{ISMN}.}
\label{fig:Rbox\runidx}
\end{figure}

\begin{figure}[ht!]
\centering
\begin{minipage}[c]{0.34\textwidth}
    \centering
    \plotfig{./qa4sm_graphics/bulk_boxplot_urmsd.png}
    \subcaption{} % (a)
\end{minipage}%
\hfill
\begin{minipage}[c]{0.64\textwidth}
    \centering
    \mapfig{./qa4sm_graphics/bulk_overview_\runpair_urmsd.png}
    \subcaption{} % (b)
\end{minipage}
\caption{Distribution of the \ac{ubRMSD} across all valid sensors (a) and spatial overview of the \ac{ubRMSD} (b) between \ac{SMOS} \ac{L2} and \ac{ISMN}.}
\label{fig:urmsd\runidx}
\end{figure}

\begin{figure}[ht!]
\centering
\begin{subfigure}{\textwidth}
    \centering
    \includegraphics[width=0.85\fullwidth,height=9.5cm,keepaspectratio]{./qa4sm_graphics/bulk_boxplot_R_metadata_lc_2010.png}
    \caption{}
\end{subfigure}
\medskip
\begin{subfigure}{\textwidth}
    \centering
    \includegraphics[width=0.85\fullwidth,height=9.5cm,keepaspectratio]{./qa4sm_graphics/bulk_boxplot_urmsd_metadata_lc_2010.png}
    \caption{}
\end{subfigure}
\caption{Pearson's R (a) and \ac{ubRMSD} (b) between \ac{SMOS} \ac{L2} and \ac{ISMN} by land cover class (\ac{ESA} \ac{CCI}, 2010) of the in situ sites. N is the number of tested sensors per class (including sensors without valid results).}
\label{fig:Rmeta\runidx}
\end{figure}

